\appendix
\section*{Numerical Differentiation}

Numerical differentiation is the process of approximating the derivative of a function using discrete values. This is essential in simulations where analytic derivatives are unknown or difficult to compute.

\subsection*{Forward Difference}

The simplest approximation is the forward difference, derived directly from the definition of the derivative:
\[
f'(x) \approx \frac{f(x+h) - f(x)}{h}
\]
Using Taylor expansion, $f(x+h) = f(x) + f'(x)h + \frac{1}{2}f''(\xi)h^2$. Substituting this reveals the truncation error:
\[
\frac{f(x+h) - f(x)}{h} = f'(x) + O(h)
\]
The forward difference is a first-order approximation.

\subsection*{Central Difference}

A more accurate approach is the central difference, which uses values on both sides of $x$:
\[
f'(x) \approx \frac{f(x+h) - f(x-h)}{2h}
\]
Consider the Taylor expansions:
\[
f(x+h) = f(x) + f'(x)h + \frac{1}{2}f''(x)h^2 + \frac{1}{6}f'''(\xi_1)h^3
\]
\[
f(x-h) = f(x) - f'(x)h + \frac{1}{2}f''(x)h^2 - \frac{1}{6}f'''(\xi_2)h^3
\]
Subtracting the two equations:
\[
f(x+h) - f(x-h) = 2f'(x)h + O(h^3)
\]
\[
\frac{f(x+h) - f(x-h)}{2h} = f'(x) + O(h^2)
\]
The central difference is a second-order approximation, meaning the error decreases quadratically as $h \to 0$.

\subsection*{Comparison}

In simulations, central difference is generally preferred over forward or backward difference because the $O(h^2)$ error term allows for much higher accuracy with the same step size $h$. However, it requires an extra function evaluation $f(x-h)$, which might not always be available (e.g., at the start of a simulation time series).
